Customer segmentation depends on the information used to describe each customer, yet representation is often treated as neutral. This dissertation compares a compact transaction representation with a nested twelve-dimensional behavioural representation for the same 11,719 RetailRocket purchasers. The cohort and evaluation design are held constant. CLASSIX supplies intrinsic aggregation provenance, while ExKMC supplies post hoc rules for K Means assignments.

Changing the representation produces markedly different partitions. Independent full-data selection gives ARI 0.014 for the two four-cluster K Means results and ARI 0.068 for CLASSIX. Across 41 CLASSIX settings acceptable in both spaces, matched-parameter ARI ranges from $-0.051$ to $0.392$. Mean fixed-parameter overlap-resampling ARI remains above 0.964, showing that a configuration can be repeatable within each representation while the representations disagree.

The explanation tradeoff also changes. At a common four-leaf budget, compact ExKMC trees reach mean held-out fidelity 0.9448 and rich trees reach 0.9650 with the same depth. At sixteen leaves the values are 0.9924 and 0.9747. Rich CLASSIX produces 559 aggregation groups compared with 131 in the compact fit. One of four rich K Means clusters passes the documented addressability audit. No compact K Means cluster or CLASSIX cluster passes.

Representation changes observable cluster structure and the conditions needed to explain it. Richer behaviour does not automatically improve internal quality or produce deployment-ready segments. Its business value remains limited to descriptive and queryable information until prospective intervention outcomes are measured.
